%!TEX root = ./../../thesis-trier.tex

\todo[inline]{hier an einigen stellen befehl semiautomat nicht benutzt, weil probleme wie unten.}

Let $\semiautomat{A}{Q}{\Sigma}{\delta}$ and $\Gamma \subseteq \Sigma$.

\begin{enumerate}
	
	\IN{restriction on automaton}
	\item The \emph{restriction of $\mathcal A$ to $\Gamma$} is $\semiautomat{A_{|\Gamma}}{Q}{\Gamma}{\delta_{|\Gamma}}$
	where $\delta_{|\Gamma} \colon Q \times \Gamma \to Q$
	is defined by $\applytrans{\delta_{|\Gamma}}{q}{x} = \applytrans{\delta}{q}{x}$ for all $q \in Q$
	and all $x \in \Gamma$.	\label{def:restriction_Gamma}
	
	\label{def:restriction_aut}
	\item Let $S \subseteq Q$ be an upper set in $\mathcal A$.
	The \emph{restriction of $\mathcal A$ to the upper set $S$}
	is the 
	deterministic and complete semiautomaton $\semiautomat{A_{|S}}{S}{\Sigma}{\delta_{|S}}$ 
	with $\delta_{|S} \colon S \times \Sigma \to S$
	given by
	$\applytrans{\delta_{|S}}{q}{x} = \applytrans{\delta}{q}{x}$
	for $q \in S$ and $x \in \Sigma$.
	We also call~$\mathcal A_{|S}$ the \emph{sub(semi-)automaton} \IN{subautomaton}
	corresponding to (or induced/defined by) the upper set $S \subseteq Q$.
	We often just denote $\delta_{|S}$ by $\delta$ in this context.
		 
	\label{def:restriction_autSGamma}
	\item If $S \subseteq Q$ is an upper set in $\mathcal A_{|\Gamma}$,
	then by %$\semiautomat{A_{|S,\Gamma}}{S}{\Gamma}{\delta_{|S,\Gamma}}$
%	$\semiautomat{{\restrictAutGamS{\mathcal A}{\Gamma}{S}}{S}{\Gamma}{\delta_{|S,\Gamma}}$ %
	$\restrictAutGamS{\mathcal A}{\Gamma}{S} = (S, \Gamma, \delta_{|\Gamma, S})$	
	we denote the \emph{restriction of $\mathcal A$ to the the alphabet $\Gamma$ and the subset $S$}, i.e.,
	$(\mathcal A_{|\Gamma})_{|S}$. 
\end{enumerate}

Upper sets can be identified with subautomata. 
Let $\semiautomat{A}{Q}{\Sigma}{\delta}$
and $S, T \subseteq Q$ be upper sets. The following facts are readily seen.
We have
\begin{equation} \label{prelim:prop:automata:isos:item:subaut}
	(\mathcal A_{|S})_{|T} \cong \mathcal A_{|T}  \text{ for } T \subseteq S
\end{equation}
Clearly, $\mathcal A_{|S}$ is strongly connected if and only if~$S$
is a sink component.

\begin{proposition}
	\label{prelim:prop:automata:isos:item:corsub}
	The extension of the inclusion map $i \colon S \to Q$
	given by $i(q) = q$ for $q \in S$ to subsets induces
	a bijection between	the strongly connected components (resp.~upper sets) in~$\mathcal A_{|S}$ 
	and the strongly connected components (resp.~upper sets) $R$ in~$\mathcal A$ with~$R \subseteq S$.
\end{proposition}